% 本文件由 LaTeX 排版 Agent 根据有效 translation.json 自动生成。
% author=Lactantius; work_id=0707; title=凤凰

\chapter{凤凰}

% structure: p0001 source=h1 target=omitted reason=page-title-already-rendered-by-parent

作者不详，但归之于\Person{拉克唐提乌斯}\par

在东方之初，有一处幸福而幽静的地方，永恒的极门敞开着。然而，它并非位于夏日或冬日初升之处，而是太阳从春之战车洒下白昼的地方。那里有一片平原，伸展着开阔的地带；没有山丘隆起，也没有幽谷凹陷。但那地方高出群山十二肘，众山之巅在我们看来已算高峻。这里是太阳的圣林；一座林木繁茂的树林，长满各种树木，因常青的荣耀而绽放。当\Person{法厄同}的火焰燃遍苍穹时，那地方未被火焰损伤；当洪水淹没世界时，它高耸于\Person{丢卡利翁}的波涛之上。无令人衰弱的疾病，无衰老之苦，无残酷的死亡，也无严酷的恐惧靠近此地；无可怕的罪恶，无疯狂的贪财之心，无\Person{玛尔斯}，也无因嗜杀而燃烧的狂怒。痛苦的悲伤不在，衣衫褴褛的匮乏、不眠的忧虑和暴烈的饥饿也都不在。没有风暴在那里肆虐，没有可怕的狂风；白霜不以其寒露覆盖大地。没有云彩在平原上伸展其羊毛般的覆盖，也没有浑浊的水汽从高处降下；但中央有一道泉水，他们称之为\Quote{活水;}它清澈、温柔，充满甘甜的水，每月一次涌出，十二次灌溉整个树林。这里有一种树，高耸的树干结出成熟的果实，不会落在地上。这座圣林，这片树林，由一只独特的鸟——\OriginalTerm{凤凰}{phoenix}——居住；它独特，却因自己的死亡而复活。它服从并侍奉\Person{福玻斯}，是一位卓越的侍者。它的本性赋予它这一职责。当藏红花的黎明初升变红，以玫瑰色的光芒驱散星辰时，它三次四次将身体浸入神圣的波涛，三次四次从活水中啜饮。它高高升起，坐在高树的最高枝头，这棵树独览整座树林；它转向初升的新生\Person{福玻斯}，等待他的光芒和初升的光辉。当太阳推开明亮大门的门槛，第一道光的微明闪耀时，它开始吟唱神圣的歌曲，以奇妙的声音迎接新光，夜莺的歌声或\Person{缪斯}的笛声都不能以\Place{库拉}之调与之相比。但据说垂死的天鹅也无法模仿它，\Person{墨丘利}的竖琴悦耳的琴弦也不能。当\Person{福玻斯}将他的马匹带回开阔的天空，并不断前行，展现他的整个圆盘时，它以三次拍翅表示赞赏，三次朝拜那载火之首后便沉默。它还能以无误的声音区分昼夜的快慢时辰：它是树林的监督者，树林的可敬女祭司，独享你的秘密，哦，\Person{福玻斯}。当它活过了一千年的寿命，岁月的漫长使它疲惫，为了更新逝去的岁月，在命运的催促下，它逃离了惯常树林的可爱的床榻。当它离开圣地，出于重生的渴望，它便寻求这个世界，那里死亡统治着。它带着满腹岁月，迅速飞向\Place{叙利亚}，\Person{维纳斯}亲自赐予那地名为\Place{腓尼基}；它穿过无路可寻的荒野，寻找那地方隐蔽的树林，那里幽谷深藏着一片遥远的森林。然后它选一棵高大的棕榈树，树梢高入云霄，因鸟而得名凤凰，那里没有有害的生物能闯入，没有粘滑的蛇，也没有猛禽。\Person{埃俄罗斯}将风关在悬挂的洞穴中，以免它们以风暴伤害明亮的空气，或以免南风收集的云朵穿过空天遮蔽太阳的光芒，阻碍这只鸟。然后它为自己建造一个巢或坟墓，因为它死去为要生存；然而它自己产生自己。于是它收集汁液和香气，\Person{亚述人}从富饶的树木中采集的，富足的\Person{阿拉伯人}采集的；\Person{俾格米各族}或\Place{印度}采摘的，或\Place{萨巴}之地从柔软胸怀产出的。于是它堆积肉桂和远香豆蔻的香气，以及混有叶子的香脂。温和的肉桂枝条和芬芳的刺檗也不缺，乳香的泪珠和丰盈的滴珠也不缺。此外，它加上繁盛的甘松的嫩穗，并加入太悦人的没药草场。随即它将要改变的身体放在铺好的巢上，安详的肢体放在这样的床榻上。然后它用嘴将汁液洒遍四周和肢体上，即将以自己的葬礼死去。于是在各种香气中它交出生命，不惧怕如此巨大托付的信实。同时，被死亡摧毁的身体，这死亡却是生命的源头，发热，这热本身产生火焰；它从远处的天空之光中接火：它燃烧，化为燃尽的灰烬。这些灰烬在死亡中收集，它仿佛使之融合成团，产生类似种子的效果。据说，从这灰烬中生出一种无肢体的动物，但那蠕虫据说是奶白色的。它以不完全成形的身体突然大大增长，并聚集成一个完美圆蛋的形状。此后它再次成形，如之前的样子，凤凰破壳而出，如田野中的毛虫，当它们被丝线系在石头上时，常能变成蝴蝶。在我们的世界里，没有食物为它预备；也没有人喂养未长羽毛的它。它啜饮从繁星之天降下的天上琼浆的微妙露珠。它收集这些；这鸟借此在香气中滋养自己，直到长成天然形状。但当它开始焕发青春时，它飞出，即将返回故地。然而在此之前，它用香脂油膏、没药和溶解的乳香包裹自己身体的一切残余、骨骸或灰烬以及自己的遗迹，并以敬虔之口将它揉成圆形，用脚携带此物，走向太阳升起之处，在祭台前停留，在神圣殿堂中取出。它向观看者展示并呈现自己，令人惊叹；如此之美，如此之荣。首先，它的颜色如同熟透的石榴在光滑果皮下所显的光彩；如同田野中罂粟所生的叶子所包含的颜色，当\Person{福罗拉}在嫣红天空下铺开她的衣裳时。它的肩膀和美丽的胸膛以此覆盖闪耀；它的头、颈和背的上部以此闪耀。它的尾巴伸展，点缀着黄色的金属，其中混合着紫色的斑点。它的双翼之间上方有一明亮的记号，如同\Person{伊里斯}在高处常给云朵染上色彩。它以绿色宝石的混合闪烁发光，纯角质的闪亮喙张开。它的眼睛大；你会相信那是两颗红宝石；从中间闪耀着明亮的火焰。一顶光环冠冕适合作蔽它的整个头部，如同高处的\Person{福玻斯}头部的荣耀。鳞片覆盖它的大腿，点缀着黄色金属，但玫瑰色以荣耀涂抹它的爪。它的形态看起来融合了孔雀与\Place{法锡斯}彩鸟的形象。\Place{阿拉伯}之地所产的有翼生物，无论是兽是鸟，几乎不能与它的大小相当。然而它不迟缓，不像那些因躯体巨大而动作笨拙、重量极重的鸟。它轻快敏捷，充满王者的美丽。它总以这样的样子出现在人眼前。\Place{埃及}人来到如此奇观面前，欢腾的人群向这稀有的鸟致敬。他们即刻在圣石上雕刻它的形象，并以新的名号标记事件和日子。各类飞鸟聚集；无一家念及猎物，无一念及恐惧。在一群鸟的陪伴下，它飞过天空，人群伴随着它，因这敬虔的职责而欢欣。但当它到达纯粹以太的区域时，它随即返回；之后它隐藏在自己的地域中。哦，幸运之鸟和命运之鸟，神自己赐予它由自身而生！无论它是雌是雄，或非男非女，或二者皆是，它是有福的，它不进入\Person{维纳斯}的盟约。死亡对它而言就是\Person{维纳斯}；它唯一的快乐在于死亡：为要出生，它渴望先前死去。它是自己的后代，自己的父亲和继承者，自己的乳母，永远是自己养育的孩子。它自己确实是，却不相同，因为它是自己，又不是自己，借着死亡的福佑获得了永生。\par

\SourceRecord{Translated by William Fletcher.；Ante-Nicene Fathers；Vol. 7.；Edited by Alexander Roberts, James Donaldson, and A. Cleveland Coxe.；Buffalo, NY: Christian Literature Publishing Co.,；1886.；<http://www.newadvent.org/fathers/0707.htm>.}
